\documentclass[a4paper,10pt]{article}

\usepackage[top=12mm,bottom=12mm,left=12mm,right=12mm]{geometry}
\usepackage{fontspec}
\usepackage{polyglossia}
\setdefaultlanguage{thai}
\setotherlanguage{english}

\setmainfont{Sarabun}[
    UprightFont = * Regular,
    BoldFont = * Bold,
    ItalicFont = * Italic,
    BoldItalicFont = * Bold Italic,
    Scale = 1.0
]
\newfontfamily\thaifont{Sarabun}[Script=Thai,
    UprightFont = * Regular,
    BoldFont = * Bold,
    ItalicFont = * Italic,
    BoldItalicFont = * Bold Italic,
    Scale = 1.0
]
\newfontfamily\thaifonttt{Sarabun}[Script=Thai,
    UprightFont = * Regular,
    BoldFont = * Bold,
    Scale = 0.95
]
\setmonofont{Sarabun}[Scale=0.95]

\usepackage{graphicx}
\usepackage{xcolor}
\usepackage{tcolorbox}
\usepackage{hyperref}

\definecolor{rbrugreen}{RGB}{16, 80, 52}
\definecolor{darknavy}{RGB}{15, 23, 42}
\definecolor{techblue}{RGB}{14, 116, 144}

\hypersetup{
    colorlinks=true,
    linkcolor=techblue,
    urlcolor=techblue,
    pdftitle={แผ่นการ์ดมาร์กเกอร์ 3D WebAR ArUco คู่มือปฏิบัติการจุลชีววิทยาการเกษตร},
    pdfauthor={ผศ.ดร.ชีวะ ทัศนา และ ผศ.ดร.จิรภัทร จันทมาลี}
}

\pagestyle{empty}

\newcommand{\arcard}[4]{%
    \begin{tcolorbox}[colback=white, colframe=gray!40, boxrule=0.6pt, arc=3pt, left=3pt, right=3pt, top=3pt, bottom=3pt, halign=center, width=\linewidth]
        \includegraphics[height=75mm,keepaspectratio]{#1}\\[1.5mm]
        {\footnotesize\bfseries\color{techblue} #2}\\[0.5mm]
        {\scriptsize\color{darknavy} #3 | #4}
    \end{tcolorbox}%
}

\begin{document}

% =============================================================
% แผ่นที่ 1: การ์ดปฏิบัติการที่ 1 - 4
% =============================================================
\begin{center}
    {\fontsize{14pt}{17pt}\bfseries\color{rbrugreen} แผ่นการ์ดมาร์กเกอร์ 3D WebAR ArUco Interactive Book}\\[1mm]
    {\fontsize{10pt}{13pt}\bfseries\color{techblue} ตำราจุลชีววิทยาการเกษตรและเทคโนโลยี AIoT (คู่มือปฏิบัติการ บทที่ 11)}\\[0.5mm]
    {\fontsize{8pt}{10pt}\color{darknavy} คณะวิทยาศาสตร์และเทคโนโลยี มหาวิทยาลัยราชภัฏรำไพพรรณี | ArUco DICT\_6X6\_250 (แผ่นที่ 1/2)}
\end{center}

\vspace{-2mm}
\begin{tcolorbox}[colback=techblue!6, colframe=techblue, boxrule=0.6pt, arc=3pt, left=6pt, right=6pt, top=3pt, bottom=3pt]
\footnotesize
\textbf{\color{techblue}คำแนะนำการใช้งาน:}\\
1. ตัดการ์ดตามกรอบ หรือส่องโดยตรงจากแผ่นพิมพ์หรือหน้าจอดิจิทัล\\
2. เปิดเบราว์เซอร์สมาร์ตโฟนเข้าสู่ WebAR: \url{http://localhost/ar_3d_aruco_book/}\\
3. ส่องกล้องให้เห็นรหัส ArUco กึ่งกลางหน้าจอ แบบจำลอง 3 มิติเชิงปฏิสัมพันธ์จะปรากฏขึ้นทันที 60 FPS
\end{tcolorbox}

\vspace{1mm}

\noindent
\begin{minipage}[t]{0.485\textwidth}
    \arcard{aruco_cards/card_lab1_aruco_0.png}{ปฏิบัติการที่ 1: Winogradsky Eco-Column (ID: 0)}{บทที่ 11}{โมเดล 3D คอลัมน์ 4 ชั้นนิเวศจุลินทรีย์}
\end{minipage}\hfill
\begin{minipage}[t]{0.485\textwidth}
    \arcard{aruco_cards/card_lab2_aruco_1.png}{ปฏิบัติการที่ 2: Rhizobium \& Leghemoglobin (ID: 1)}{บทที่ 11}{ปมรากถั่วผ่าซีก และโครงสร้างโปรตีน 1GDJ}
\end{minipage}

\vspace{2mm}

\noindent
\begin{minipage}[t]{0.485\textwidth}
    \arcard{aruco_cards/card_lab3_aruco_2.png}{ปฏิบัติการที่ 3: Phosphate Solubilizer (PSI) (ID: 2)}{บทที่ 11}{วุ้น Pikovskaya วงใส ละลายหินฟอสเฟต}
\end{minipage}\hfill
\begin{minipage}[t]{0.485\textwidth}
    \arcard{aruco_cards/card_lab4_aruco_3.png}{ปฏิบัติการที่ 4: Trichoderma Confrontation (ID: 3)}{บทที่ 11}{Dual Culture PIRG และเอนไซม์ Chitinase}
\end{minipage}

\clearpage

% =============================================================
% แผ่นที่ 2: การ์ดปฏิบัติการที่ 5 - 8
% =============================================================
\begin{center}
    {\fontsize{14pt}{17pt}\bfseries\color{rbrugreen} แผ่นการ์ดมาร์กเกอร์ 3D WebAR ArUco Interactive Book}\\[1mm]
    {\fontsize{10pt}{13pt}\bfseries\color{techblue} ตำราจุลชีววิทยาการเกษตรและเทคโนโลยี AIoT (คู่มือปฏิบัติการ บทที่ 11)}\\[0.5mm]
    {\fontsize{8pt}{10pt}\color{darknavy} คณะวิทยาศาสตร์และเทคโนโลยี มหาวิทยาลัยราชภัฏรำไพพรรณี | ArUco DICT\_6X6\_250 (แผ่นที่ 2/2)}
\end{center}

\vspace{-2mm}
\begin{tcolorbox}[colback=rbrugreen!6, colframe=rbrugreen, boxrule=0.6pt, arc=3pt, left=6pt, right=6pt, top=3pt, bottom=3pt]
\footnotesize
\textbf{\color{rbrugreen}ระบบติดตามตำแหน่งเชิงแสง 6 มิติ (6-DoF Optical Tracking):}\\
รหัส ArUco แต่ละแผ่นเข้ารหัสเมทริกซ์ 6$\times$6 ป้องกันการหมุนหลงทิศแบบ 360 องศา คำนวณเวกเตอร์ตำแหน่ง $\vec{t}$ และควอเทอร์เนียน $\mathbf{q}$ รองรับการหมุนรอบวัตถุเสมือนจริง 60 FPS บนเว็บเบราว์เซอร์ทุกแพลตฟอร์ม
\end{tcolorbox}

\vspace{1mm}

\noindent
\begin{minipage}[t]{0.485\textwidth}
    \arcard{aruco_cards/card_lab5_aruco_4.png}{ปฏิบัติการที่ 5: AMF Root Colonization (ID: 4)}{บทที่ 11}{รากพืชย้อมสี Trypan Blue และกริดนับสปอร์}
\end{minipage}\hfill
\begin{minipage}[t]{0.485\textwidth}
    \arcard{aruco_cards/card_lab6_aruco_5.png}{ปฏิบัติการที่ 6: Biofertilizer QC \& Coliform (ID: 5)}{บทที่ 11}{หลอดเจือจาง $10^{-1}$ ถึง $10^{-6}$ และ MacConkey}
\end{minipage}

\vspace{2mm}

\noindent
\begin{minipage}[t]{0.485\textwidth}
    \arcard{aruco_cards/card_lab7_aruco_6.png}{ปฏิบัติการที่ 7: Soil Metagenomic DNA (ID: 6)}{บทที่ 11}{สกัดหลอดบีดบีตติ้ง เกลียวคู่ B-DNA 1BNA}
\end{minipage}\hfill
\begin{minipage}[t]{0.485\textwidth}
    \arcard{aruco_cards/card_lab8_aruco_7.png}{ปฏิบัติการที่ 8: PCR Diagnosis \& Vision AI (ID: 7)}{บทที่ 11}{เครื่อง PCR เทอร์โมไซเคิล เจล UV และ AI}
\end{minipage}

\end{document}
